\begin{longtable}{|p{3.5cm}|p{6cm}|p{3cm}|}
\captionsetup{justification=raggedright, singlelinecheck=false}
\caption{Rencana Implementasi Modul \textit{Exfiltration}} \label{tbl:rencana-exfiltration} \\
\hline
\textbf{Langkah} & \textbf{Penjelasan} & \textbf{Tools} \\
\hline
\endfirsthead

\hline
\textbf{Langkah} & \textbf{Penjelasan} & \textbf{Tools} \\
\hline
\endhead

\hline
\endfoot

Menentukan server penerima &
Menyusun layanan sederhana pada sisi penerima untuk menerima data hasil eksfiltrasi dari host Windows. &
Python \textit{FastAPI} \\
\hline

Mempersiapkan struktur data yang akan dikirim &
Data hasil \textit{crawling} yang sudah disimpan sementara dalam \textit{buffer in memory} dikemas menjadi format JSON agar mudah dikirim dan tidak menimbulkan \textit{burst traffic} mencolok. &
\textit{Modul} formatting JSON \\
\hline

Pengaturan interval pengiriman acak (\textit{low and slow}) &
Setiap \textit{payload} dikirim dalam interval waktu acak dan tidak beraturan sehingga pola komunikasi tidak menunjukkan ritme repetitif yang dapat memicu deteksi anomali. &
Timer/Scheduler \\
\hline

Mengirimkan data melalui kanal komunikasi umum (HTTP/S) &
Eksfiltrasi dilakukan menggunakan \textit{request} HTTP/HTTPS lokal sehingga pola komunikasinya menyerupai aplikasi biasa. &
PowerShell \\
\hline

Pembersihan \textit{buffer} setelah pengiriman &
Setelah setiap potongan data dikirim, \textit{buffer} dihapus dari memori untuk meningkatkan \textit{stealth} dan mengurangi jejak aktivitas pada sistem. &
Memory buffer module \\
\hline

Verifikasi pengiriman data &
Memastikan setiap data berhasil diterima oleh server lokal tanpa \textit{error}. &
Log server \\
\hline

\end{longtable}